

\chapter{机械波}

根据底层物理机制与相互作用的不同，宏观与微观世界中的波动现象主要分为以下三类：

\begin{definition}[机械波，电磁波，物质波]
机械波：机械振动的状态以一定的速度在弹性介质中由近及远地传播出去过程。
\begin{itemize}
    \item 产生条件一（波源）：机械振动的产生者，负责提供系统最初的扰动能量。
    \item 产生条件二（弹性介质） ：传播振动的承担者（如水、空气、绳索等）。由于机械波依赖质点间的弹性力传递扰动，故真空中无法传播机械波（例如真空中无法传声）。
\end{itemize}
电磁波：变化的电场和变化的磁场在空间中相互激发、交替产生并由近及远传播的过程。电磁波的传播在本质上不需要任何依赖实体介质，其可以在真空中高速自由传播（如光波、无线电波、X射线等）。

物质波（也称概率波）是微观粒子（如电子、质子、中子等）本身所具有的一种固有属性。 物质波具有与宏观机械波完完全全不同的物理性质。它在宏观上并不代表实体的连续波动，而是描述粒子在空间某处出现的概率分布，其运动规律严格遵从量子力学理论。
\end{definition}

\begin{definition}[波阵面、波前和波线、波长、周期、频率和波速]
\begin{itemize}
\item 波阵面（波面）：在波动传播过程中，空间中具有相同振动相位的各质点所连成的连续曲面。
\item 波前：在任何特定时刻，传播在最前方（最外边缘）的那一个波阵面。
\item 波线：沿波的传播方向所作的一系列带箭头的几何线，其指向表示波在该处的能量传播方向。在各向同性的均匀介质中，波线与波阵面在几何上总是保持相互垂直。
\end{itemize}
为了定量描述波动在时间和空间双重维度的交织特征，我们定义以下四个核心特征物理量：
\begin{itemize}
\item 波长 $\lambda$：沿波的传播方向，两个相邻的、振动状态（相位）完全相同的振动质点之间的空间距离。它在物理上代表一个完整波形的几何长度。
\item 周期 $T$：波前进一个波长的距离所需要的时间。它在数值上等于介质中任一质点完成一次完全往复振动所需的时间。
\item 频率 $\nu$：单位时间内波动所传播的完整波的数目。它等于周期的倒数（$\nu = 1/T$）。
\item 波速 $u$：在波动过程中，某一振动状态（或相位）单位时间内所传播的距离。
\end{itemize}
波的周期 $T$ 或频率 $\nu$ 只决定于波源的振动。当一列波从一种介质跨越边界进入另一种介质时，其频率 $\nu$ 保持绝对不变。
$$T = \frac{\lambda}{u}, \qquad \nu = \frac{1}{T} = \frac{u}{\lambda},~~\boxed{u = \frac{\lambda}{T} = \lambda \nu}$$
\end{definition}

\begin{definition}[应力、应变与弹性模量]\label{def:elastic-modulus}
材料力学中描述"形变有多大"和"恢复力有多强"的基本量：
\begin{itemize}
 \item \textbf{应力} \(\sigma=F/S\)：单位面积上的内力，单位 \unit{Pa}。想象把一根杆拉长——杆内部任意截面上，两侧材料互相拉扯的力除以截面积就是应力。
 \item \textbf{应变} \(\varepsilon=\Delta l/l\)：相对形变量，无量纲。一根 \SI{1}{m} 的杆被拉长了 \SI{1}{mm}，应变就是 \(10^{-3}\)。
 \item \textbf{弹性模量}：应力与应变之比，衡量"材料有多硬"。
 \begin{itemize}
 \item 杨氏模量 \(Y=\sigma/\varepsilon\)：描述拉伸/压缩的刚度。钢的 \(Y\approx\SI{200}{GPa}\)，橡胶的 \(Y\) 只有几 MPa。
 \item 切变模量 \(G\)：描述剪切形变（平行于截面的滑动）的刚度。
 \item 体积模量 \(B=-V\dd p/\dd V\)：描述均匀压缩时体积变化的刚度。水的 \(B\approx\SI{2.2}{GPa}\)。
 \end{itemize}
\end{itemize}
模量越大，材料越"硬"——同样的力只能产生更小的形变。后面推导波的能量时会再次用到 \(\sigma=Y\varepsilon\)。
\end{definition}
决定波速的核心有两方面：
\begin{itemize}
 \item 介质的\textbf{弹性}有多强，也就是它偏离平衡后有多想恢复；
 \item 介质的\textbf{惯性}有多大，也就是它有多不容易被带动。
\end{itemize}
弹性越强，回复越快，波就更容易传得快；密度越大，单位体积介质越“沉”，越不容易被带动，波速就会变慢。先看分子对应什么模量、分母对应什么密度。模量越大，恢复越快；密度越大，惯性越大，波速就越慢。
\begin{itemize}
 \item 固体中纵波：\(u=\sqrt{\frac{Y}{\rho}}\)，其中 \(Y\) 是杨氏模量，对应伸缩形变，\(\rho\) 是体密度。
 \item 固体中横波：\(u=\sqrt{\frac{G}{\rho}}\)，其中 \(G\) 是切变模量，对应剪切形变。
 \item 液体和气体中的纵波：\(u=\sqrt{\frac{B}{\rho}}\)，其中 \(B\) 是体积模量，对应容积形变。
 \item 绳中横波：\(u=\sqrt{\frac{T}{\mu}}\)，其中 \(T\) 是张力，\(\mu\) 是质量线密度。这里不用体弹性模量，是因为弦被横向拨开后真正负责“拉回去”的不是体形变，而是原先绷紧的张力。
\end{itemize}
